\begin{proof}
Reemplazando $A$ por un anillo de fracciones $A_f$ de modo que
$D(f)\subset\operatorname{Reg}(\operatorname{Spec}(A))$, se puede suponer ya
que el anillo $A$ es \emph{regular}. Por hipótesis existe en $K'$ un sistema
$(t_i)_{1\leq i\leq n}$ de elementos algebraicamente independientes sobre
$K$ tales que $K'$ es una extensión algebraica separable finita de
$K(t_1,\ldots,t_n)$; considerando un sistema finito
\oldpage[IV-2]{165}
de generadores $t'_j$ $(1\leq j\leq m)$ de $K'$ sobre
$K(t_1,\ldots,t_n)$, que se pueden suponer enteros sobre
$A_1=A[t_1,\ldots,t_n]$, se ve que
$A'_1=A_1[t'_1,\ldots,t'_m]$ es \emph{finita} sobre $A_1$ y tiene por cuerpo de
fracciones $K'$. Si se pone $X'=\operatorname{Spec}(A')$ y
$X'_1=\operatorname{Spec}(A'_1)$, del hecho de que los cuerpos de funciones
racionales $R(X')$ y $R(X'_1)$ son ambos isomorfos a $K'$, y de que $X'$
y $X'_1$ son $A$-preesquemas de tipo finito, se deduce que existen un
abierto $U'\subset X'$ y un abierto $U'_1\subset X'_1$ que son
$A$-\emph{isomorfos} (\textbf I, 6.5.5). La cuestión se reduce, por tanto, a demostrar que
$\operatorname{Reg}(X'_1)$ contiene un abierto no vacío; en otras palabras,
se puede suponer que $A'$ es una $A_1$-\emph{álgebra finita}. Ahora bien, se sabe
(\textbf 0, 17.3.7) que $A_1$ es un anillo \emph{regular}, por lo que es posible limitarse al
caso en que $A'$ es una $A$-álgebra \emph{finita} y $K'$ una extensión \emph{separable y
finita} de $K$. Si $\xi$ es el punto genérico de $X$, $A'_\xi=K'$ es
entonces un módulo libre sobre $A_\xi=K$; por tanto (Bourbaki,
\emph{Alg. comm.}, cap. II, \S5, n\textsuperscript{o} 1,
cor. de la prop. 2), reemplazando si es necesario $A$ por un $A_f$,
se puede suponer que $A'$ es un $A$-módulo \emph{libre} de tipo finito. Sea entonces
$(x_h)_{1\leq h\leq r}$ una base de este $A$-módulo, y pongamos
\[
d=\det(\operatorname{Tr}_{A'/A}(x_hx_k))=
\det(\operatorname{Tr}_{K'/K}(x_hx_k))\in A.
\]
Como $K'$ es separable sobre $K$, se sabe (Bourbaki, \emph{Alg.}, cap. IX, \S2,
prop. 5) que $d\neq0$; reemplazando, si es necesario, $A$ por el anillo de
fracciones $A_d$, se puede suponer que $d$ es \emph{invertible} en $A$. Pero entonces,
para todo $z\in\operatorname{Spec}(A)$, si se denota por $\bar x_h$
$(1\leq h\leq r)$ la imagen canónica de $x_h$ en
$A'(z)=A'\otimes_A \boldsymbol{k}(z)$, se tiene
$\det(\operatorname{Tr}_{A'(z)/\mathbf{k}(z)}(\bar x_h\bar x_k))=\bar d$, imagen
canónica de $d$ en $\boldsymbol{k}(z)=A_z/\mathfrak m_z$, y como $\bar d$ es invertible
(y, por tanto, $\neq0$) en $\boldsymbol{k}(z)$, se sabe \emph{(loc. cit.)} que $A'(z)$ es una
$\boldsymbol{k}(z)$-álgebra \emph{separable}, es decir, un producto directo de cuerpos que son
extensiones finitas y separables de $\boldsymbol{k}(z)$. Como tal álgebra es un anillo
\emph{regular}, se ve que el morfismo $g:X'\to X$ es \emph{plano} y que sus fibras
$g^{-1}(z)$ son \emph{regulares} para todo $z\in X$; se concluye entonces de
(6.5.2, (ii)) que $X'$ es \emph{regular}, lo que termina la demostración en este
caso.
\end{proof}
